% BEGIN VOL07-EXPANSION VII07
\section{Supplementary worked examples}
\label{sec:vii07-pedagogy}

\begin{example}[A trivial bundle and its sections]\label{ex:vii07-ped-01}
The projection \(\pi:M\times\mathbb R^k\to M\) is the trivial rank-\(k\) bundle.  A smooth section has the form \(s(p)=(p,\sigma(p))\) for a smooth map \(\sigma:M\to\mathbb R^k\).  Thus ordinary vector-valued functions are exactly sections of a trivial vector bundle.
\end{example}

\begin{example}[The tangent bundle of the circle is trivial]\label{ex:vii07-ped-02}
On \(S^1\subset\mathbb R^2\), the vector field \(V(x,y)=(-y,x)\) is smooth and nowhere zero.  The map
\[
S^1\times\mathbb R\longrightarrow TS^1,\qquad ((x,y),a)\longmapsto aV(x,y)
\]
is a vector-bundle isomorphism, so \(TS^1\) is a trivial line bundle.
\end{example}

\begin{example}[The Möbius line bundle]\label{ex:vii07-ped-03}
Start with \([0,1]\times\mathbb R\) and identify \((0,v)\sim(1,-v)\).  The resulting bundle over \(S^1\) is locally a product but has a sign change in its transition function.  This is the standard example showing that local triviality does not imply global triviality.
\end{example}

\section{Graded supplementary exercises}
The following exercises add a deliberate progression from concrete calculation to structural proof, hypothesis testing, application, and synthesis.

\subsection*{Standard computations and constructions}

\begin{exercise}[Section of a trivial bundle]\label{exr:vii07-ped-01}
Describe all sections of \(M\times\mathbb R\to M\).
\end{exercise}
\begin{hint}
A section must project to the identity.
\end{hint}
\begin{solution}
They are exactly maps \(s(p)=(p,f(p))\) with \(f\in C^\infty(M)\).
\end{solution}

\begin{exercise}[Circle tangent section]\label{exr:vii07-ped-02}
Verify that \(V(x,y)=(-y,x)\) is tangent to \(S^1\).
\end{exercise}
\begin{hint}
Take its dot product with the radius vector.
\end{hint}
\begin{solution}
\((x,y)\cdot(-y,x)=0\), so \(V\) lies in the tangent line at every point.
\end{solution}

\begin{exercise}[Rank of a Whitney sum]\label{exr:vii07-ped-03}
If \(E\) has rank \(r\) and \(F\) rank \(s\), find the rank of \(E\oplus F\).
\end{exercise}
\begin{hint}
Look fiberwise.
\end{hint}
\begin{solution}
Each fiber is \(E_p\oplus F_p\), of dimension \(r+s\).
\end{solution}

\begin{exercise}[Transition matrix]\label{exr:vii07-ped-04}
For two trivializations of a rank-\(k\) bundle, what kind of map is the transition function?
\end{exercise}
\begin{hint}
Compare the two linear coordinates on each fiber.
\end{hint}
\begin{solution}
It is a smooth map from the overlap to \(GL(k,\mathbb R)\).
\end{solution}

\begin{exercise}[Pullback fiber]\label{exr:vii07-ped-05}
For \(f:N\to M\) and a bundle \(E\to M\), identify the fiber of \(f^*E\) over \(q\in N\).
\end{exercise}
\begin{hint}
Use the definition of the pullback bundle.
\end{hint}
\begin{solution}
\((f^*E)_q\cong E_{f(q)}\).
\end{solution}

\subsection*{Proofs}

\begin{exercise}[Zero section is smooth]\label{exr:vii07-ped-06}
Prove that the zero section of any vector bundle is smooth.
\end{exercise}
\begin{hint}
Check it in a local trivialization.
\end{hint}
\begin{solution}
In local bundle coordinates it is \(p\mapsto(p,0)\), a smooth map.
\end{solution}

\begin{exercise}[Transition cocycle]\label{exr:vii07-ped-07}
Prove that transition functions satisfy \(g_{ij}g_{jk}=g_{ik}\) on triple overlaps.
\end{exercise}
\begin{hint}
Compose the three changes of fiber coordinates.
\end{hint}
\begin{solution}
Changing from chart \(k\) to \(j\) and then \(j\) to \(i\) is the same change as going directly from \(k\) to \(i\), yielding the cocycle equation.
\end{solution}

\begin{exercise}[Pullback of a trivial bundle]\label{exr:vii07-ped-08}
Prove that the pullback of a trivial rank-\(k\) bundle is trivial.
\end{exercise}
\begin{hint}
Write the pullback explicitly.
\end{hint}
\begin{solution}
If \(E=M\times\mathbb R^k\), then \(f^*E=\{(q,(f(q),v))\}\cong N\times\mathbb R^k\) by \((q,(f(q),v))\mapsto(q,v)\).
\end{solution}

\begin{exercise}[Nowhere-zero section trivializes a line bundle]\label{exr:vii07-ped-09}
Prove that a real line bundle with a nowhere-zero smooth section is trivial.
\end{exercise}
\begin{hint}
Use the section as a basis in each one-dimensional fiber.
\end{hint}
\begin{solution}
The map \(M\times\mathbb R\to L\), \((p,a)\mapsto a\,s(p)\), is smooth and fiberwise linear with smooth local inverse, hence a bundle isomorphism.
\end{solution}

\subsection*{Counterexamples and hypothesis tests}

\begin{exercise}[Local triviality implies global triviality]\label{exr:vii07-ped-10}
Test this claim with the Möbius line bundle.
\end{exercise}
\begin{hint}
Track the sign change around the base circle.
\end{hint}
\begin{solution}
False.  The Möbius bundle is locally a product but has nontrivial monodromy and no nowhere-zero section that is compatible with a global product coordinate.
\end{solution}

\begin{exercise}[Every vector bundle has a nowhere-zero section]\label{exr:vii07-ped-11}
Test the claim with \(TS^2\).
\end{exercise}
\begin{hint}
Use the hairy-ball theorem.
\end{hint}
\begin{solution}
False.  The tangent bundle of the two-sphere has no nowhere-zero continuous section.
\end{solution}

\begin{exercise}[Bundle rank may jump]\label{exr:vii07-ped-12}
Test whether a vector bundle can have different fiber dimensions on one connected base.
\end{exercise}
\begin{hint}
Use local triviality.
\end{hint}
\begin{solution}
False.  Rank is locally constant because a local trivialization fixes the fiber dimension, and hence is constant on connected components.
\end{solution}

\subsection*{Applications and investigations}

\begin{exercise}[Tangent bundle as velocities]\label{exr:vii07-ped-13}
Explain why \(TM\) is the natural state space for positions and instantaneous velocities.
\end{exercise}
\begin{hint}
A point of \(TM\) is a pair \((p,v)\) with \(v\in T_pM\).
\end{hint}
\begin{solution}
The bundle projection remembers position, while the fiber records all possible velocities at that position.
\end{solution}

\begin{exercise}[Bundle-valued fields]\label{exr:vii07-ped-14}
Interpret a smooth vector field as a section of a bundle.
\end{exercise}
\begin{hint}
Use the tangent bundle projection.
\end{hint}
\begin{solution}
A vector field chooses smoothly one vector \(X_p\in T_pM\) at each point, exactly a smooth section \(X:M\to TM\).
\end{solution}

\subsection*{Challenge problems}

\begin{exercise}[Möbius obstruction]\label{exr:vii07-ped-15}
Show directly that the Möbius line bundle has no nowhere-zero section.
\end{exercise}
\begin{hint}
Represent a section by a function \(f:[0,1]\to\mathbb R\) with \(f(1)=-f(0)\).
\end{hint}
\begin{solution}
If \(f(0)\neq0\), the endpoint condition gives opposite signs, so the intermediate value theorem forces a zero.  If \(f(0)=0\), the section already vanishes.
\end{solution}

\begin{exercise}[Bundle isomorphism test]\label{exr:vii07-ped-16}
Suppose \(F:E\to E'\) covers the identity on \(M\) and is a linear isomorphism on every fiber.  Explain why a smooth \(F\) is a vector-bundle isomorphism.
\end{exercise}
\begin{hint}
Work in local trivializations and invert the matrix-valued fiber map.
\end{hint}
\begin{solution}
Locally \(F(p,v)=(p,A(p)v)\) with \(A(p)\in GL(k)\).  Matrix inversion is smooth on \(GL(k)\), so the local inverses patch to a smooth bundle inverse.
\end{solution}

% END VOL07-EXPANSION VII07
